% Epilogue of Singular Learning Theory for the Rest of Us (written by the coordinating author from
% the chapters' final ledger cards).
\specialchapter{Epilogue}{The ledger}{{0/2,2/0}{1}{1}{a^2+b^2}}
\sltfloatprefix{E}
\begin{refsection}[chapters/16-epilogue.bib]

The prologue fitted the line $y=abx$ to pure noise and found that its two parameters learn like one: the free energy grows like $\tfrac12\log n-\log\log n$, not like $\log n$. Fifteen chapters later the same replacement has been made in the formulas of Bayesian statistics that charge a model $\tfrac d2$ for its $d$ parameters. Each chapter recorded its result on a card with the same six fields: the quantity, its value in the regular theory, its value in the singular theory, a running example, what was measured, and what the result rests on. Read one at a time, the cards summarize theorems. Read side by side, they show what the theory is: a geometric number and its multiplicity, which govern most of what a Bayesian asks of a model, under hypotheses that deep learning does not always meet. This epilogue sets the cards side by side, sets the hypotheses against practice, and collects the open problems.

\addsec{The ledger}

\Cref{tab:epilogue:ledger1,tab:epilogue:ledger2} collect the cards of the fifteen chapters, with each field compressed to a phrase, the running example folded into the singular column, and what the results rest on gathered in \cref{tab:epilogue:assumptions}; the precise statements, with their hypotheses, are in the chapters. Read down the regular column, and one number recurs: $\tfrac d2$, which is half a unit of $\log n$ per parameter in the free energy, $\tfrac1{2n}$ per parameter in the generalization error, and $\tfrac1{2\beta}$ per parameter in the energy of a tempered posterior. Read down the singular column, and the learning coefficient $\lambda$ takes its place in each of these roles, with its multiplicity $m$ beside it.

\begin{table}[p]
\begin{fullwidth}
\footnotesize
\renewcommand{\arraystretch}{1.18}
\begin{tabularx}{\linewidth}{@{}>{\raggedright\arraybackslash}p{25mm}>{\hsize=0.72\hsize\raggedright\arraybackslash}X>{\hsize=1.2\hsize\raggedright\arraybackslash}X>{\hsize=1.08\hsize\raggedright\arraybackslash}X@{}}
\toprule
\textsc{\color{singular}chapter: quantity} & \textsc{\color{regular}regular theory} & \textsc{\color{singular}singular theory} & \textsc{\color{singular}measured}\\
\midrule
\multicolumn{4}{@{}l}{\textsc{\color{singular}part i. where regular theory ends}}\\[0.3ex]
1\enspace regular models: $F_n-nS_n$ and $\E[G_n]-S$ & $\frac d2\log n+O_p(1)$ and $\frac d{2n}+o(\frac1n)$: Laplace's method, BIC and AIC & the special case $\lambda=\frac d2$, $m=1$ & slopes $0.499$, $0.998$, $0.992$ against $\log n$ for $d=1,2,2$; $n$ times Laplace's error on the mixture at a regular truth $0.596$ at $n=10^4$, against $0.594$ predicted \\
2\enspace singular models: the true set and the shape of the posterior & $W_0$ a point, the Fisher information invertible, the posterior close to a Gaussian of width $n^{-1/2}$ & $W_0$ contains curves on which the Fisher information is singular; for $y=abx$, $\log\abs a/\log n\to\Unif(-\frac12,0)$ under the posterior; no Gaussian limit in $C^1$ coordinates & posterior quartiles of $\log\abs a/\log n$ at $n=10^{16}$: $-0.3800$, $-0.2530$, $-0.1265$, against $-0.3809$, $-0.2539$, $-0.1270$ predicted \\
3\enspace one coefficient: $F_n-nS_n$ when the likelihood depends on $w$ through $c=g(w)$ & $\frac d2\log n$: $\log n$ for $y=abx$, $\frac32\log n$ for $y=a_1a_2a_3x$ & $\lambda\log n-(m-1)\log\log n+O_p(1)$, with $\lambda$ and $m$ read off the prior density of $c$ at the truth, and the limit law of the remainder; $y=abx$: $\lambda=\frac12$, $m=2$ & slopes $0.41$, $0.32$, $0.23$, $0.16$ over $10^2\le n\le10^8$ for $\lambda=\frac12$ ($m=2$), $\frac12$ ($m=3$), $\frac14$, $\frac16$ \\[0.4ex]
\multicolumn{4}{@{}l}{\textsc{\color{singular}part ii. the learning coefficient}}\\[0.3ex]
4\enspace volume: the prior mass $V(\eps)$ of $\{K<\eps\}$ & $C\eps^{d/2}$, the volume of an ellipsoid & $C\eps^\lambda(\log\frac1\eps)^{m-1}$, with $-\lambda$ the largest pole of the zeta function and $m$ its order; the quartic cup: $V(\eps)=1.311\,\eps^{1/2}$ exactly for $\eps\le\frac12$, against $\eps^1$ & Monte Carlo exponents $1.0006$, $0.5001$, $0.7494$, $0.6665$ for $\lambda=1,\frac12,\frac34,\frac23$; for the product model $0.5001$ with $\eps^{1/2}(c_0+c_1\log\frac1\eps)$ fitted, $0.39$ for a straight line \\
5\enspace blow-ups: $\lambda$ and $m$ on a resolution & $\frac d2$ and $m=1$: by the Morse lemma, one blow-up resolves a nondegenerate minimum & $\min_j(h_j+1)/(2k_j)$ over the charts, $m$ the most ties at one point; $\frac12(a^2+b^4)$: two blow-ups, $\lambda=\frac34$; $(a^2-b^3)^2$: three, $\frac5{12}$ & chart sums equal the original volumes to $10^{-6}$; local slopes $0.7500$, $0.4134$, $0.7500$ at $\eps\approx10^{-14}$ for $\lambda=\frac34,\frac5{12},\frac34$ \\
6\enspace Newton polygons: $\lambda$ from the exponents of $K$ & every exponent that matters is $2$: $\frac d2$, $m=1$ & $\lambda=1/t$, where the diagonal meets the Newton polygon at $(t,t)$, and $m=2$ if it meets a vertex (proved for sums of even monomials with positive coefficients in two variables, stated for nondegenerate losses); $\lambda\le1/t$ in every coordinate system; $\frac12(a^4+b^6)$: $\frac5{12}$ & for thirty random sums, quadrature exponents within $3\times10^{-7}$ of $1/t$; Monte Carlo slopes at $\eps\ge10^{-4}$ low by $0.09$ when $m=2$ \\
7\enspace the catalogue: $\lambda$ and $m$ of models in use & $\frac d2$, or half the dimension of the set of distributions & rank-one regression at the zero truth: $\frac12\min(M,N)$, with $m=2$ if and only if $M=N$; reduced-rank regression: the formula of Aoyagi and Watanabe, and deep linear networks: Aoyagi's, both stated and checked; $M=N=3$: $\frac32$ at the truth $0$, against $3$ & two-temperature estimates at $n=10^5$: $1.40$, $2.50$, $2.50$, $3.44$ for $\lambda=\frac32,\frac52,\frac52,\frac72$, below $\lambda$ where $m=2$ by somewhat less than the predicted correction \\
\bottomrule
\end{tabularx}
\end{fullwidth}
\caption{The ledger, Parts I and II. Each row compresses the card of its chapter; the chapters state the hypotheses. Here $d$ is the number of parameters, $F_n-nS_n$ the normalized free energy, $\E[G_n]-S$ the Bayes generalization error, $K=L-L(w_0)$ the excess population loss, which vanishes on the true set $W_0$, $V(\eps)$ the prior volume of $\{K<\eps\}$, and $M$ and $N$ the numbers of outputs and inputs of rank-one regression; measured slopes are fitted against $\log n$ or $\log\frac1\eps$.}\label{tab:epilogue:ledger1}
\end{table}

\begin{table}[p]
\begin{fullwidth}
\footnotesize
\renewcommand{\arraystretch}{1.18}
\begin{tabularx}{\linewidth}{@{}>{\raggedright\arraybackslash}p{25mm}>{\hsize=0.72\hsize\raggedright\arraybackslash}X>{\hsize=1.2\hsize\raggedright\arraybackslash}X>{\hsize=1.08\hsize\raggedright\arraybackslash}X@{}}
\toprule
\textsc{\color{singular}chapter: quantity} & \textsc{\color{regular}regular theory} & \textsc{\color{singular}singular theory} & \textsc{\color{singular}measured}\\
\midrule
\multicolumn{4}{@{}l}{\textsc{\color{singular}part iii. watanabe's theorems}}\\[0.3ex]
8\enspace free energy: $F_n-nS_n$ & $\frac d2\log n+O_p(1)$: Schwarz's BIC & $\lambda\log n-(m-1)\log\log n+O_p(1)$: proved for regressions $y=G(w)x+\varepsilon$, Watanabe's theorem for analytic models; $y=ax_1+b^3x_2$: $\frac23$ & slopes $0.998$, $0.749$, $0.667$, $0.491$ for $\lambda=1,\frac34,\frac23,\frac12$, and $0.394$ for the product model; the singular BIC finds the true rank in $182$ of $200$ data sets at $n=300$, BIC in $5$ \\
9\enspace generalization: $\E[G_n]-S$ & $\frac d{2n}+o(\frac1n)$; training error $-\frac d{2n}$ & $\frac\lambda n+o(\frac1n)$: proved for one-coefficient regressions with Gaussian inputs and noise, Watanabe's theorem in general; by his theorem, training error $(\lambda-2\nu)/n$; $y=ab^3x$: $\frac1{6n}$ & $n(\E[G_n]-S)$ at $n=10^6$: $0.164\pm0.024$, $0.244\pm0.031$, $0.504\pm0.047$ for $\lambda=\frac16,\frac14,\frac12$; $0.422\pm0.038$ for the product model, which approaches $\frac12$ like $\frac12-1/\log n$ \\
10\enspace WAIC: estimating $G_n$ from one posterior & AIC, $L_n(\hat w)+\frac dn$ or $T_n+\frac dn$, unbiased to $o(\frac1n)$; leave-one-out agrees & $\E[\WAIC]=\E[G_n]+o(\frac1n)$ (Watanabe); leave-one-out $=\WAIC+O_p(n^{-2})$, proved for $y=abx$; penalty $V_n/n$ with $\E[V_n]\to2\nu$; at the singular truth of $y=abx$, $2\nu=1$, half of AIC's $2$ & bias of WAIC zero within $\pm0.13/n$ for $10^2\le n\le10^5$; bias of $T_n+2/n$ between $1.02/n$ and $1.19/n$ \\
11\enspace temperature: the tempered mean energy, WBIC, $\hat\lambda$ & equipartition, $\frac d{2\beta}$; $\WBIC=\BIC+o_p(1)$ & population mean energy $\frac\lambda\beta-\frac{m-1}{\beta\log n\beta}+\cdots$; $F_n-\WBIC=-(m-1)\log\log n+O_p(\sqrt{\log n})$: proved for realizable linear models, Watanabe's theorem in general; the quartic model: $\frac3{4\beta}$ & $\hat\lambda$ at $n=10^8$: $0.750$, $0.667$, $1.000$ for $\lambda=\frac34,\frac23,1$; $0.438$ for the product model, its population value \\
12\enspace phases: the posterior mass of a region & every region pays $\frac d2\log n$; the posterior goes to the region of the best fit & local free energies $nL_n(w_\alpha)+\lambda_\alpha\log n-(m_\alpha-1)\log\log n$, proved for linear models, and the smallest wins; transitions in $n$ lower the loss and raise $\lambda$; the double well: from $\lambda=\frac12$ to $1$ near $n=3000$ & posterior mass of the regular region $0.24$, $0.53$, $0.97$ at $n=10^3$, $3162$, $10^4$, against $0.25$, $0.54$, $0.96$ predicted and $0.75$, $0.93$, $1.00$ from the regular theory \\[0.4ex]
\multicolumn{4}{@{}l}{\textsc{\color{singular}part iv. deep learning}}\\[0.3ex]
13\enspace local learning coefficient: $\lambda(\wstar)$ at a minimum, and its estimate & Laplace at $\wstar$: $\frac d2$, or $\frac r2$ for a Hessian of rank $r$ & $\frac r2\le\lambda(\wstar)\le\frac r2+\frac{d-r}4$, equal to $\frac r2$ exactly when the true set is a manifold of dimension $d-r$ near $\wstar$; rank-one regression with $M=N=3$: $\frac32$ at $0$ and at $(e_1,0)$ & the exact estimator for the ReLU neuron at $n=10^4$: $0.32$ ($m=2$) and $0.46$ ($m=1$) for $\lambda=\frac12$; SGLD errors from $-82\%$ to $+118\%$ across step sizes; a fixed Gaussian localizer gives $0.99$ for $\frac32$ when $N=2$ \\
14\enspace stages: $\lambda$ at the critical points that training visits & $\frac d2=9$ at every nondegenerate minimum of the toy model of superposition & $\frac{d-s}2$, $m=1$, at a minimally singular critical set of dimension $s$: $7$ for the $5$-gon, $\frac{17}2$ for the $5^+$-gon and the $6$-gon, computer-assisted; the order of stages conjectured from the posterior's & at the exact polygons $7.06$, $8.34$, $8.73$ at $n=10^6$, but $8.21$, $8.32$, $9.63$ at $n=10^3$; along SGD, $4.58\to7.56$ and $5.20\to7.83$; $9$ of $32$ runs dropped, none to the $6$-gon \\
15\enspace SGD is not Bayes: where training stops and where the posterior goes & the same place: the minimizer is the center of an asymptotically normal posterior & on $y=abx$, gradient flow conserves $a^2-b^2$: $\log\abs a/\log n\to\frac12\delta_{-1/2}+\frac12\delta_0$ from the prior, against $\Unif(-\frac12,0)$ under the posterior & at $n=10^{12}$, mass of the middle half of the scales $0.490$ under the posterior, $0.005$ under the flow from the prior; at $n=10^4$, $36\%$ of SGD runs collapse onto the origin \\
\bottomrule
\end{tabularx}
\end{fullwidth}
\caption{The ledger, Parts III and IV. Here $\nu$ is the singular fluctuation, $T_n$ the Bayes training loss, $V_n$ the functional variance, $\beta$ the inverse temperature, $\hat\lambda$ the two-temperature estimate of $\lambda$, $L_n(w_\alpha)$ the empirical loss at the best point of a region, $\wstar$ a local minimum, and $s$ the dimension of a set of critical points.}\label{tab:epilogue:ledger2}
\end{table}

\addsec{What the ledger shows}

\paragraph{One number in four roles.} The regular theory does not distinguish the four places where $\tfrac d2$ appears, because there they are equal. The singular theory separates them and then proves that one number fills all four: the exponent of the volume of nearly true parameters (\cref{ch:volume}), the coefficient of $\log n$ in the free energy (\cref{ch:coefficient,ch:freeenergy}), $n$ times the expected Bayes generalization error (\cref{ch:generalization}), and the slope of the tempered posterior's mean energy in $1/\beta$ (\cref{ch:temperature}). The chapters measured it in all four roles, by procedures that share no computation: Monte Carlo volumes of the prior, quadrature of the marginal likelihood over many data sets, the divergence from the truth to the predictive distribution, and posterior energies at two temperatures. On the running models with $m=1$ the four measurements agree with each other and with the value that Part II computes by changes of coordinates, and with $m\ge2$ they agree once the correction described next is made. That agreement is the theory's main empirical content, and it is why one number can serve as a singular model's complexity. The second number of Part III is not $\lambda$: the singular fluctuation $\nu$, which $\lambda$ does not determine, enters the training error and the Gibbs errors and sets the penalty of WAIC (\cref{ch:generalization,ch:waic}). For the product model at its singular truth $2\nu=1=2\lambda$, but for the cusp $y=ab^2x$, $2\nu\approx0.70$ while $2\lambda=\frac12$, and the chapters know of no closed form for $\nu$, even for the power-law class of one-coefficient models (\cref{ch:generalization}). WAIC needs neither number, which is what makes it useful.

\paragraph{The multiplicity is slow.} Wherever $m\ge2$, the measurements fall short of $\lambda$: the free-energy slope $0.39$ of the product model over $10^2\le n\le10^6$ (the prologue and \cref{ch:freeenergy}), its generalization error $0.422/n$ at $n=10^6$ (\cref{ch:generalization}), its two-temperature estimate $0.438$ at $n=10^8$ (\cref{ch:temperature}), the rank-one estimates of \cref{ch:catalogue}, and the local estimate $0.32$ at the origin of the ReLU neuron against $0.46$ on one of its arms (\cref{ch:llc}). The shortfall is of the order of $(m-1)/\log n$, and where the chapters computed the correction, the corrected prediction fits the measurement, or nearly; none of these is a failure of the theory. But a correction of order $1/\log n$ does not go away at the sample sizes met in practice: at $n=10^{12}$ it is still $0.036$ for $m=2$. An experimenter who reads $\lambda$ off a slope without knowing $m$ reads it low (\cref{ch:temperature,ch:llc}), and the chapters know of no estimator of $m$ with a guarantee (\cref{ch:coefficient,ch:volume}). Of what the headline formula leaves out, the multiplicity is the part that biases every measurement of $\lambda$ in the same direction; the bounded terms are the part that moves the transitions of the posterior at practical sample sizes (\cref{ch:phases,ch:stages}).

\paragraph{Where the posterior goes.} Three rows replace the regular picture of the posterior, a Gaussian at the maximum-likelihood estimate. The posterior of $y=abx$ spreads over scales along its true set instead of concentrating at the rate $n^{-1/2}$ (\cref{ch:singular}); when several regions nearly fit the data, it moves between them as $n$ grows, trading loss against learning coefficient (\cref{ch:phases}); and gradient flow, which does stop at a point, stops at a point chosen by its dynamics, through a conserved quantity, not by volume (\cref{ch:sgd}). In a regular model the optimizer and the posterior end up in the same place, and the theory of the one is a theory of the other. In a singular model they find the same function and part ways about the parameters. This is why Part IV measures the learning coefficient at a trained network instead of predicting it: the network is not a draw from the posterior, and Part III does not say which point training reaches.

\paragraph{What is proved and what is measured.} The certainty of the ledger decreases down the page. Parts I and II prove their central results completely, or reduce them to Hironaka's theorem, stated in \cref{ch:toolkit}; what they take from the literature, such as the generalization error of a general regular model (\cref{ch:regular}), the theorems on Newton polyhedra beyond positive sums in two variables (\cref{ch:newton}), and the formulas of Aoyagi and Watanabe (\cref{ch:catalogue}), they state with its source and check numerically where they can. Part III proves its theorems completely for classes of models in which the computation can be carried through, one coefficient (\cref{ch:generalization}), the product model itself (\cref{ch:waic}), or regressions linear in the inputs (\cref{ch:freeenergy,ch:temperature,ch:phases}), with one limit, the Bayes training error, only argued (\cref{ch:generalization}), and states Watanabe's general theorems with their sources and an account of what the general proofs add. In Part IV the theorems concern the geometry of a minimum, the Hessian bounds (\cref{ch:llc}) and the coefficient of a minimally singular critical set (\cref{ch:stages}), and the dynamics of gradient flow on $y=abx$ (\cref{ch:sgd}). Under Watanabe's conditions the estimator of the local learning coefficient is consistent for exact posteriors with a compact localizer, by his theorem on WBIC applied to a ball; what SGLD computes on a network has no guarantee, and the stages of training are measurements with a conjectured Bayesian explanation.

\addsec{The hypotheses against practice}

Every card ends with what its result rests on. \Cref{tab:epilogue:assumptions} collects the hypotheses that recur, sets each against how networks are trained, and says what the book shows without it. Two of them fail for every large network in use: the number of parameters is comparable to $n$ or larger, not fixed while $n$ grows, and the parameters come from an optimizer, not from the posterior. The first is the deeper gap, because every asymptotic statement of the book is a statement about $n\to\infty$ with $d$ fixed; the second is the one that \cref{ch:sgd} makes concrete.

\begin{table}[tp]
\begin{fullwidth}
\footnotesize
\renewcommand{\arraystretch}{1.18}
\begin{tabularx}{\linewidth}{@{}>{\raggedright\arraybackslash}p{30mm}>{\raggedright\arraybackslash}p{24mm}>{\raggedright\arraybackslash}X>{\raggedright\arraybackslash}X@{}}
\toprule
\textsc{\color{singular}hypothesis} & \textsc{\color{singular}used in} & \textsc{\color{regular}in deep learning} & \textsc{\color{singular}without it}\\
\midrule
the model is fixed as $n\to\infty$ & every asymptotic theorem & $d$ is comparable to $n$ or larger, and models grow with the data & for linear models with many fewer inputs than samples the free energy is pinned down only to $O_p(k)$, for $k$ coefficients (\cref{ch:freeenergy}); bounds for redundant networks of fixed size, such as that of \textcite{NagayasuWatanabe2023} for deep ReLU networks, are a start (\cref{ch:sgd}); otherwise open \\
realizability, or relatively finite variance & \cref{ch:coefficient,ch:freeenergy,ch:generalization,ch:waic} & no network is the true distribution & Watanabe's theorems need only relatively finite variance and a unique optimal distribution; without the first, in his non-renormalizable example, the generalization error is of order $n^{-2/3}$ by his sketch, the free energy has a term of order $n^{1/3}$ by the book's argument, and WAIC stays unbiased in simulation (\cref{ch:freeenergy,ch:waic}) \\
the loss is analytic near the true set & \cref{ch:blowups,ch:freeenergy,ch:generalization,ch:waic} & ReLU networks are only piecewise analytic & $\lambda$ and $m$ keep their meaning as exponents of the volume where it has power-logarithm bounds (\cref{ch:blowups}), which nothing in the book guarantees for a ReLU network (\cref{ch:volume}); the Hessian bounds need only a $C^4$ loss (\cref{ch:llc}); there is no resolution theorem to compute $\lambda$ \\
a compact parameter set; a prior positive at the true set & \cref{ch:regular,ch:freeenergy,ch:llc} & Gaussian initialization, weight decay, unbounded weights & in rank-one regression with $N<M$, a fixed Gaussian localizer makes the population estimate converge to $\frac N2$, the value at a more singular point, not to $\lambda(\wstar)=\frac M2$ (\cref{prop:llc:gaussian}): the localization must depend on $n$ \\
exact posterior averages & \crefrange{ch:singular}{ch:phases} & parameters from SGD; averages over short SGLD chains & gradient flow selects its end point by a conservation law and SGD by drifting to balance, not by volume (\cref{ch:sgd}); SGLD estimates have no guarantee and move with the step size (\cref{ch:llc}) \\
independent data from a fixed distribution & every chapter with data & augmentation, curricula, data that shift during training & nothing in this book, beyond a cited remark that with fixed or dependent inputs WAIC, unlike leave-one-out, still estimates the conditional generalization loss (\cref{ch:waic}) \\
transitions in $n$ & \cref{ch:phases} & stages are seen in training time & that stages of training have Bayesian antecedents is a conjecture (\cref{conj:stages:bah,conj:sgd:bah}) \\
asymptotics in $n$, with corrections in $1/\log n$ & \cref{ch:coefficient,ch:generalization,ch:temperature} & $n$ is finite, and $\log\log n$ is not small next to $\log n$ & the multiplicity biases estimates of $\lambda$ by about $(m-1)/\log n$ (\cref{ch:catalogue,ch:temperature,ch:llc}); constants of local free energies decide transitions at practical $n$ (\cref{ch:phases,ch:stages}) \\
\bottomrule
\end{tabularx}
\end{fullwidth}
\caption{The hypotheses of the book against the practice of deep learning.}\label{tab:epilogue:assumptions}
\end{table}

\addsec{Open problems}

The chapters end with open problems, and they fall into five groups.

\paragraph{Computing $\lambda$.} Exact learning coefficients are known for rank-one and reduced-rank regression, for deep linear networks, for one tanh layer with a single input, and for a few other families, among them some mixtures, naive Bayesian networks, and latent trees; they are not known for one tanh layer with several inputs and outputs and more than four hidden units, for two nonlinear layers beyond the smallest, or for hidden Markov models or attention layers, and the exact ReLU values claimed by \textcite{Aoyagi2025} await independent checks (\cref{ch:catalogue}). No practical algorithm is known for the real log canonical threshold of a loss in many variables (\cref{ch:blowups}). In three or more variables \textcite[Section~1]{Greenblatt2010} asks when coordinates exist in which the Newton polyhedron gives sharp estimates; for learning, the question is for which architectures and truths a short, explicit sequence of blow-ups makes the loss nondegenerate in every chart (\cref{ch:newton}). For losses that are not analytic, such as those of ReLU networks, a theory of learning coefficients would need a replacement for Hironaka's theorem (\cref{ch:blowups}).

\paragraph{Estimating $\lambda$ and $m$.} No estimator of the volume function of an analytic loss in high dimension is known to have controlled relative error at a cost polynomial in $d$ and $\log\frac1\eps$ (\cref{ch:volume}), and the chapters found no estimator of $\lambda$ with a proved error of smaller order than $1/\log n$ when $m\ge2$ is unknown: the analyses of \textcite{Imai2019a} and \textcite{TakioSuzuki2025} give consistency at the rate $(\log n)^{-1/2}$ and do not treat the bias $-(m-1)/\log n$ (\cref{ch:temperature}). For the local learning coefficient, the conjecture of \textcite[Appendix~D]{LauEtAl2025} asks for conditions on the localizer, which must depend on $n$, under which the estimator converges to $\lambda(\wstar)$ (\cref{ch:llc}); why stochastic-gradient MCMC chains recover local volume exponents at all is posed by \textcite{HitchcockHoogland2025}; and no method is known that gives confidence intervals covering the bias of an estimate, in a large network or along training (\cref{ch:llc,ch:stages}).

\paragraph{Finite $n$ and growing $d$.} Almost every error term in Part III is asymptotic. Non-asymptotic bounds for the generalization error (\cref{ch:generalization}) and for WBIC and $\hat\lambda$ (\cref{ch:temperature}) are open, and so is the behavior of the free energy when $d$ grows with $n$ (\cref{ch:freeenergy}). The lower-order terms of local free energies, which decide transitions at practical sample sizes, are unknown in general, and \textcite{CarrollEtAl2025} list them among the open questions of their account (\cref{ch:coefficient,ch:phases,ch:stages}); so is the coefficient of $1/\log n$ in the mean free energy of a model with several coefficients and $m\ge2$ (\cref{ch:coefficient}).

\paragraph{The posterior and training.} Beyond one coefficient, for the tanh neuron, the mixture, and networks at reducible truths, the law of the posterior over scales is unknown (\cref{ch:singular}). That the transitions met in training have Bayesian antecedents, transitions of the posterior in $n$ between the same regions, is the hypothesis of \textcite{ChenEtAl2023}, and what links the two is open (\cref{ch:phases,ch:stages,ch:sgd}); even for $y=abx$ the long-run law of stochastic gradient descent with a constant step, and its relation to the posterior's local free energies, are open (\cref{ch:sgd}). And why natural data and architectures give low loss, learning coefficient, and singular fluctuation is the question of \textcite{Murfet2023a} (\cref{ch:sgd}).

\paragraph{Watanabe's conjectures.} Watanabe conjectured that the singular fluctuation is governed by a generalized Fisher information on the set of true parameters \parencite[Section~5]{Watanabe2010a} (\cref{ch:generalization}), and that the parity on which the error of WBIC depends is a birational invariant for every truth \parencite[p.~875]{Watanabe2013} (\cref{ch:temperature}). Whether WAIC agrees with leave-one-out cross-validation to order $n^{-2}$ for models whose loss is smooth but not analytic, and what replaces its theorem without relatively finite variance, are open as well (\cref{ch:waic}).

\addsec{What to read next}

Watanabe's books are the sources of the general theory. \Textcite{Watanabe2009} proves the free energy formula and the equations of state with resolution of singularities and empirical processes, and this book's appendix states what it uses from there; the unrealizable case, WBIC, and the equivalence of WAIC with cross-validation came later, in papers that the chapters cite. \Textcite{Watanabe2018} is a graduate text with exercises, on WAIC, WBIC, and cross-validation among much else; the exercise books of \textcite{Suzuki2023a,Suzuki2023b} work through the same theory with code; and \textcite{Watanabe2024} reviews the theory and its recent advances. For computing learning coefficients, the thesis of \textcite{Lin2011} develops the algebra behind Part II, and \textcite{DrtonPlummer2017}, with its published discussion, is the statistician's way into model selection with the singular BIC. For deep learning, \textcite{WeiEtAl2023} argue for deep-learning readers that singularity is a feature of Bayesian learning, not a defect; \textcite{LauEtAl2025} introduce the local learning coefficient; \textcite{HooglandEtAl2025} use it to locate stages of development in transformers; and the sequence of \textcite{Carroll2023} is an informal introduction for readers from machine learning.

\addsec{Two parameters, again}

The prologue's model is the book in miniature. Its parameters fold onto the line of slopes, so its true set is a cross (\cref{ch:singular}). Its free energy pays $\tfrac12\log n-\log\log n$ because the region of nearly true parameters is a thin cross whose arms thicken near the center (\cref{ch:coefficient,ch:volume}). Its loss is already a monomial, so the computations of Part II take one line, and the two exponents tie, which is where its multiplicity $m=2$ comes from (\cref{ch:blowups,ch:newton}). Its generalization error approaches $\frac1{2n}$ only like $(\frac12-\frac1{\log n})\frac1n$ (\cref{ch:generalization}); the penalty of WAIC is half of AIC's (\cref{ch:waic}); and its posterior spreads over scales while gradient descent picks one (\cref{ch:singular,ch:sgd}). A network has many such folds, and they meet.

Singular learning theory does not explain why networks trained by stochastic gradient descent generalize, and \cref{ch:sgd} shows on the smallest example how training and the posterior part ways. What it does establish is narrower and firmer. In Bayesian learning the complexity of a model at a truth is a geometric number, not a count of parameters; the number can be computed in small models and estimated in large ones, though there the estimates have no guarantee and their biases are known only from small models; and the classical formulas that count parameters are the special case in which the geometry is trivial. Whether the networks we train live in the regime where that geometry decides what they learn is a question about networks, not about the theory. The theory makes it possible to ask it precisely.

\notesandrefs

The cards summarize the chapters, and the references are there. The tables quote the chapters' measured numbers as the scripts printed them; the chapters give their intervals and the number of data sets behind each.

\printchapterbib
\end{refsection}
\sltfloatrestore
